% !PRINTED 516
\weberpage{516}

These results are of great importance and deserve to be
specially emphasised.

Following \gesperrt{Kronecker}, if $F(x) = 0$ is the given
equation and $N$ the associated \gesperrt{Galois} field, the quantities of
$N$ are called \gesperrt{the natural irrationalities of the equation
$F(x) = 0$}.

We then have the theorem:

\gesperrt{Every possible reduction of the Galois group is
brought about by the adjunction of a natural irrationality. If
$j$ is the index of the reduced group with respect to the
original one, then the reduction cannot come about through the
adjunction of a field of degree lower than $j$. But the degree
can be equal to $j$, and then the adjoined irrationality is a
natural one. If the degree of the adjoined field is greater
than $j$, then it is a multiple of $j$.}

\begin{center}
\S.\,158.

\gesperrt{Imprimitive Groups.}
\end{center}

We make application of our general theorems to groups of a
special nature. First we take the \gesperrt{imprimitive groups}. We have
seen in \S.\,151 that, if $f(x) = 0$ is an irreducible equation
of degree $n$ with imprimitive group, then $n$ can be decomposed
into two factors $r \boldsymbol{.} s$ in such a way that the
roots split into $s$ rows of $r$ terms each:

\weq{1}{\begin{aligned}
A &= \alpha,\ \alpha_1 \dots \alpha_{r-1},\\
B &= \beta,\ \beta_1 \dots \beta_{r-1},\\
&\cdots\cdots\cdots\cdots\cdots\cdots\cdots\cdots\\
S &= \sigma,\ \sigma_1 \dots \sigma_{r-1},
\end{aligned}}

and that by the permutations of the group $P$ the quantities of
the individual rows are not separated from one another, but are
only permuted among themselves, while on the other hand the
rows in turn are permuted among one another.

Among the permutations of the \gesperrt{Galois} group $P$ of $f(x)$
there will now be certain ones, whose totality we shall denote by $Q$,
%CONT>>
